\section{Background}
\label{sec:background}

\paragraph{The model.}
Qwen3.8-27B interleaves Gated DeltaNet layers, whose state is a fixed-size
matrix per head, with full-attention layers, whose KV cache grows with the context. Only the
attention layers pay per-token memory, which is why a 27B model can hold a 256K-token context in
a few gigabytes. The recurrent state, however, cannot simply be truncated when a draft token is
rejected: engines must either snapshot it or replay the accepted tokens.

\paragraph{Speculative decoding.}
A drafter proposes $k$ tokens; the target model scores all of them in one pass and keeps the
longest prefix that matches its own choice (greedy) or passes a rejection test (sampling), so
the output distribution is unchanged. The speed-up depends on the
mean number of tokens accepted per pass and on how cheap the drafter and the $(k{+}1)$-token
verification are. Qwen3.8-27B ships an MTP head that drafts one token per
step from the target's last hidden state. DFlash2 is a separate small
block-diffusion model that drafts a whole block in one pass, conditioned on the target's hidden
states, and is accepted further on predictable text such as code.

\paragraph{KV-cache quantization.}
The KV cache can be stored in fewer bits than the activations that produce it.
 Keys are more sensitive than values, since an error in a key changes
every attention weight it takes part in. Rotations spread outliers across channels before
quantization, and lattice codes such as $E_8$
quantize eight values jointly at a lower error than scalar rounding at the same rate.
Table~\ref{tab:kvtypes} lists the formats compared in this work.

\begin{table}[t]
  \centering
  \caption{KV-cache formats compared. Bits are per element, scales excluded.}
  \label{tab:kvtypes}
  \resizebox{\columnwidth}{!}{%
  \begin{tabular}{@{}llccl@{}}
    \toprule
    Engine & Format & K bits & V bits & Encoding \\
    \midrule
    llama.cpp & \texttt{q8\_0}/\texttt{q8\_0} & 8 & 8 & blocks of 32, one scale \\
    llama.cpp & \texttt{q8\_0}/\texttt{turbo4} & 8 & 4 & TurboQuant values (baseline) \\
    OpenJet & \texttt{q8} & 8 & 8 & as \texttt{q8\_0} \\
    NInfer & \texttt{int8} & 8 & 8 & groups of 64 \\
    NInfer & \texttt{rk8v4} & 8 & 4 & rotated keys, groups of 64 \\
    NInfer & \texttt{rk4v4-e8} & 4 & 4 & rotated, $E_8$ lattice keys \\
    \bottomrule
  \end{tabular}}
\end{table}
